\documentclass{article}
\usepackage{import}
\usepackage{tcolorbox}
\usepackage{amsmath}
\usepackage{amssymb}
\usepackage{amsthm}
\usepackage{geometry}
\usepackage{enumitem}
\usepackage{pdfpages}
\usepackage{transparent}
\usepackage{booktabs}
\usepackage{bm}
\usepackage{tabularx}
\usepackage{multirow}
% \usepackage{minted}
\usepackage{xcolor}
\usepackage{setspace}
\usepackage[hidelinks]{hyperref}
\usepackage{pgfplots}
% \usepackage{soulpos}

\tcbuselibrary{theorems}
\tcbuselibrary{skins}

\setcounter{tocdepth}{6}
\setcounter{secnumdepth}{6}

\newcommand{\param}{\boldsymbol{\theta}}
\newcommand{\bs}{\boldsymbol}
\newcommand{\data}{\mathcal D}

% pgfplots settings
\pgfplotsset{compat=1.18}
\usepgfplotslibrary{external}

% tabularx Settings
\renewcommand\tabularxcolumn[1]{>{\centering\arraybackslash}m{#1}}

% Page Settings
\newgeometry{vmargin={15mm}, hmargin={15mm, 15mm}}
\setlength{\parindent}{0pt}
\setlength{\parskip}{0.7\baselineskip}%{0.8em plus 1pt minus 1pt}

\onehalfspacing

% hyperref Settings
\hypersetup{
    colorlinks=true,
    linkcolor=blue,
    filecolor=magenta,      
    urlcolor=cyan,
}

% inkscape Settings
\newcommand{\incfig}[1]{%
    \def\svgwidth{\columnwidth}
    \import{./figure/}{#1.pdf_tex}
}

% inline code block settings
\newtcbox{\mybox}[1][]{
    on line,
    arc=1pt, outer arc=2pt,
    colback=lightgray!50!white, colframe=lightgray!50!white,
    boxsep=0pt, left=0pt, right=-1pt, top=2pt, bottom=0pt,
    boxrule=0pt, #1
}

\NewDocumentCommand{\code}{v}{%
  \begingroup\ttfamily
  \codeulinner{#1}%
  \endgroup%
}

% Definition Box Styles
\newtcbtheorem[auto counter, number within=subsection]{defn}{Definition}{
% basic settings
    theorem name and number,
    separator sign={},
    description delimiters={(}{)},
    enhanced jigsaw,
    sharp corners,
% box margins and rules
    boxrule=1pt,
    left=0.2cm,right=0.2cm,top=0.2cm,
    toptitle=0.1cm,
    bottomtitle=0.08cm,
% box colors
    colframe=white!25!black,
    colback=white,
    coltitle=white,
% box fonts
    fonttitle=\bfseries,fontupper=\normalsize
}{defn}

% Theorem Box Styles
\newtcbtheorem[auto counter, number within=subsection]{thm}{Theorem}{
% basic settings
    theorem name and number,
    separator sign={},
    description delimiters={(}{)},
    enhanced jigsaw,
    sharp corners,
% box margins and rules
    boxrule=1pt,
    left=0.2cm,right=0.2cm,top=0.2cm,
    toptitle=0.1cm,
    bottomtitle=0.08cm,
% box colors
    colframe=white!25!black,
    colback=white,
    coltitle=white,
% box fonts
    fonttitle=\bfseries,fontupper=\normalsize
}{thm}

% Lemma Box Styles
\newtcbtheorem[number within=tcb@cnt@thm]{lem}{Lemma}{
% basic settings
    theorem name and number,
    separator sign={},
    description delimiters={(}{)},
    enhanced jigsaw,
    sharp corners,
% box margins and rules
    boxrule=1pt,
    left=0.2cm,right=0.2cm,top=0.2cm,
    toptitle=0.1cm,
    bottomtitle=0.08cm,
% box colors
    colframe=white!50!black,
    colback=white,
    coltitle=white,
% box fonts
    fonttitle=\bfseries,fontupper=\normalsize
}{lem}

% Corollary Box Styles
\newtcbtheorem[number within=tcb@cnt@thm]{cor}{Corollary}{
% basic settings    
    theorem name and number,
    separator sign={},
    description delimiters={(}{)},
    enhanced jigsaw,
    sharp corners,
% box margins and rules
    boxrule=1pt,
    left=0.2cm,right=0.2cm,top=0.2cm,
    toptitle=0.1cm,
    bottomtitle=0.08cm,
% box colors
    colframe=white!50!black,
    colback=white,
    coltitle=white,
% box fonts
    fonttitle=\bfseries,fontupper=\normalsize
}{cor}

% Remark Box Styles
\newtcbtheorem[number within=subsection]{rem}{Remark}{
% basic settings
    theorem name and number,
    separator sign={},
    description delimiters={(}{)},
    enhanced jigsaw,
    sharp corners,
% box colors
    colframe=white!60!black,
    colback=white,
    coltitle=black,
% box margins and rules
    boxrule=1pt,
    left=0.2cm,right=0.2cm,top=0.2cm,
    toptitle=0.1cm,
    bottomtitle=0.08cm,
% box fonts
    fonttitle=\bfseries,fontupper=\normalsize
}{rem}

% Example Box Styles
% Remark Box Styles
\newtcbtheorem[number within=subsection]{exm}{Example}{
% basic settings
    theorem name and number,
    separator sign={},
    description delimiters={(}{)},
    enhanced jigsaw,
    sharp corners,
% box colors
    colframe=white!60!black,
    colback=white,
    coltitle=black,
% box margins and rules
    boxrule=1pt,
    left=0.2cm,right=0.2cm,top=0.2cm,
    toptitle=0.1cm,
    bottomtitle=0.08cm,
% box fonts
    fonttitle=\bfseries,fontupper=\normalsize
}{exm}


% Proof Box Styles

\title{Regularization Methods}
\author{Hyoung Joon, Kim}
\date{\today}

\begin{document}

\maketitle
\tableofcontents
\newpage

\section{Introduction}
    In this document, we discuss about methods of regularization.
    Regularization is used to deal with overfitting. There are several methods
    and these methods can be used together.

\section{Main Ideas}
    \subsection{Inductive Bias}
        \textbf{Inductive bias} is a preference for one choice over others. It is also called prior knowledge, literally.
        Inductive bias helps to constrain the space of solutions. To exploit the inductive bias, we use the following regularization approaches:

    \subsection{Weight Decay}
        \textbf{Weight decay} is used to control the size of the parameters.
        It prevents parameters being too large or too small.
        It is called weight decay, due to the following equation of \textbf{sum-of-squares regularizer}:
        \[
        \nabla E(\textbf{w}) + \lambda\textbf{w}.
        \]
        The sequential update will lead weights to zero.

        \subsubsection{Generalized Weight Decay}
        \begin{defn}{Generalized Weight Decay}{GenWeightDecay}
        A generalized weight decay regularizer is defined as follows:
        \[
        \Omega(\textbf{w}) = \frac{\lambda}{2}\sum_{j=1}^{M}\left|w_j\right|^q
        \]
        \end{defn}
        $q=2$ corresponds to the \textbf{L2 regularizer}, and $q=1$ corresponds to the \textbf{L1 regularizer}.
        The regression with L2 regularizer is called \textbf{ridge regression}, and the one with L1 regularizer is called \textbf{lasso regression}.
        The lasso regression makes \emph{sparse model}, since the lasso regression drives the coefficients or weights, to zero.

        In the Bayesian perspective, ridge regression and lasso regression is same as setting the prior as
        Gaussian and Laplace with Gaussian liklihood, respectively.

        \textbf{Elastic net} is regularizer which combines L2 and L1 regularizer;
        \[
        \Omega(\textbf{w}) = \frac{\lambda_1}{2}\sum_{j=1}^{M}\left|w_j\right|
        + \frac{\lambda_2}{2}\sum_{j=1}^{M}\left|w_j\right|^2
        \]

        \subsubsection{Constant Regularizers}
        The L2 regularizer and other regularizers, cannot include the inductive bias such as invariance and equivariance in itself.
        To deal with it, we must treat the weights and biases differently at each deep network layer. This method, which preserves the invariance is called
        \textbf{constant regularizer}. But this method is not used so much.

    \subsection{Early Stopping}
        \textbf{Early Stopping}, is to stop the training at the point of smallest error with respect to validation error.
        Stopping a training before a minimum of the training error has been reached can be understood as limiting the effective network complextity.
        
        \subsubsection{Double Descent}
            However, there is a phenomenon called \textbf{double descent} (Belkin et al., 2019). The validation error decreases first,
            and increases a bit, then suddenly, it decreases again. This can be understood as the transition of regime;
            \begin{enumerate}
                \item Under-parameterized Regime: The complexity of the model is smaller than the number of samples.
                    The model follows classical bias-variance tradeoff.
                \item Over-parameterized Regime: The complexity of the model is large enough so that the test error is 0.
                    In this case, the model's test error decreases as the model's complexity increases.
            \end{enumerate}
            For more information, check Nakkiran et al., 2019. This paper tries to explain the double descent with \textbf{effective model complexity},
            defined as following:
            \begin{defn}{Effective Model Complexity}{effModelComp}
                The \textbf{Effective Model Complexity} (EMC) of a trainging procedure $\mathcal{T}$,
                with respect to distribution $\mathcal{D}$ and parameter $\epsilon > 0$ is defined as:
                \[
                \textbf{EMC}_{\mathcal{D},\epsilon}(\mathcal{T})=
                \max{\{n|\mathbb{E}_{S\sim\mathcal{D}^n}[\text{Error}_S(\mathcal{T}(S))] \leq \epsilon\}}
                \]
                where Error$_S(M)$ is the mean error of model $M$ on train samples $S$.
            \end{defn}
            (Cannot understand the meaning of it now...)

        
    \subsection{Parameter Sharing}
        \textbf{Parameter Sharing} is to group the parameters and let the grouped parameters to have the same value.
        Parameter sharing is used to encode some invariance or equivariance into a network. CNN is the example of the parameter sharing.

        \subsubsection{Soft Parameter Sharing}
            Instead of fixing the parameters to share same values, we can encourage the parameters in the same group to have
            similar values, leading to \textbf{soft parameter sharing}. We can use \emph{mixture of Gaussains} to achieve this.

    \subsection{Dropout}
        \textbf{Dropout} is to delete the nodes from the network randomly during training. This method can be understood as implicit model averaging.
        Dropout can be applied to the input and hidden layers, but not the output layers. The probability $p=0.5$ of deleting the node, from the experiences,
        is known to work well with hidden layers. For the input layer, $p=0.8$ is works well.
        
        The dropout should be applied only when the model is being trained. However, there is a method called \textbf{Monte Carlo dropout}, where we
        still use dropout while inferencing. The Monte Carlo dropout simply make inferences mutiple times and average the results. Averaging 10 to 20 results
        is known to process good results.
        Monte Carlo dropout can provide us the uncertainty of the result, since we can check the standard deviation of the model's prediction.

\section{The sections of this note may be created to individual documents later}

\end{document}